\documentclass[12pt,a4paper]{article}
\usepackage[utf8]{inputenc}
\usepackage[french]{babel}
\usepackage[T1]{fontenc}
\usepackage{amsmath}
\usepackage{amsfonts}
\usepackage{amssymb}
\usepackage{graphicx}
\usepackage{hyperref}
\usepackage{enumitem}
\usepackage{geometry}
\usepackage{tcolorbox}

\geometry{margin=2.5cm}

\title{Guide de Validation et de Test d'Extraction\\ \large - GitLab KPI Dashboard}
\author{}
\date{} % Supprime la date automatique

\begin{document}

\maketitle

\begin{tcolorbox}[colback=blue!5!white,colframe=blue!75!black,title=Objectif du Document]
Ce guide détaille les étapes séquentielles pour réaliser une extraction complète des données GitLab. Le respect de l'ordre des étapes est important pour garantir l'intégrité des indicateurs de performance.
\end{tcolorbox}

\section{Phase 1 : Initialisation du Système}

\subsection{Étape 1 : Configuration de la Connexion GitLab}
Avant toute opération, assurez-vous que le pont entre l'application et GitLab est actif :
\begin{itemize}[label=\textbullet]
    \item Menu : \textbf{Configuration GitLab}.
    \item Action : Vérifiez que l'URL de base et le Token d'accès sont valides.
\end{itemize}

\subsection{Étape 2 : Définition des Périodes d'Analyse}
\begin{itemize}[label=\textbullet]
    \item Menu : \textbf{Gestion des Périodes} (dans la section administration sous configuration).
    \item Action : Créez ou sélectionnez la période (Mois/Année) de test.
\end{itemize}

\section{Phase 2 : Importation et Cartographie}

\subsection{Étape 3 : Importation Massive des Développeurs}
\begin{enumerate}
    \item Menu : \textbf{Import Développeurs}.
    \item Fichier : Sélectionnez \texttt{real\_senior\_devs.csv}.
    \item \textbf{Configuration Critique :}
    \begin{itemize}[label=\textbf{!}]
        \item Désactivez l'option \textbf{"Dry Run"}.
        \item Cochez impérativement les trois options : \textit{Assign to Sites}, \textit{Assign to Groups}, \textit{Assign to Projects}.
    \end{itemize}
    \item Validation : Cliquez sur \textbf{Lancer l'import}.
\end{enumerate}

\subsection{Étape 4 : Vérification des Business Units}
\begin{itemize}[label=\textbullet]
    \item Menu : \textbf{Business Units}. Vérifiez la répartition des effectifs par site.
\end{itemize}

\section{Phase 3 : Cycle d'Extraction et Visualisation}

\subsection{Étape 5 : Lancement du Moteur d'Extraction}
\begin{itemize}[label=\textbullet]
    \item Menu : \textbf{Moteur d'Extraction}. Lancez la collecte sur les projets cibles.
\end{itemize}

\subsection{Étape 6 : Monitoring et Visualisation}
\begin{itemize}[label=\textbullet]
    \item Menu : \textbf{Lots d'Extraction} pour le statut.
    \item Menus : \textbf{Commit Hub} et \textbf{Merge Request Hub} pour visualiser les données brutes collectées.
\end{itemize}

\vspace{1cm}

\begin{tcolorbox}[colback=orange!10!white,colframe=orange!75!black,title=⚠️ TRAVAIL EN COURS : MOTEUR KPI]
\textbf{Note importante :} La partie concernant les indicateurs de performance avancés (Calculs des taux MR/Commits par dev, graphiques analytiques... selon fichier spec) est actuellement en cours de développement. Ces fonctionnalités seront activées et intégrées dans la version de livraison prochaine.
\end{tcolorbox}

\end{document}
